\section{Problem Formulation}
\label{sec:formulation}
\subsection{Fixed-Horizon Rest-to-Rest Task}
Consider $N$ labeled UAVs in a known workspace $\mathcal W=[\underline{\mathbf w},\overline{\mathbf w}]\subset\R^3$. Static obstacles are closed axis-aligned bounding boxes (AABBs), $\mathcal O_b=[\underline{\mathbf o}_b,\overline{\mathbf o}_b]$, $b=1,\ldots,B$. Each UAV is represented by a sphere of radius $r$, with a required surface clearance $c$. Task identities are fixed. UAV $i$ moves from its assigned start $\mathbf s_i$ to goal $\mathbf g_i$ over the prescribed common flight duration $T_f>0$. This duration is fixed before planning and shared by every UAV; it is distinct from computation time. The boundary conditions are
\begin{equation}
 \bp_i(0)=\mathbf s_i,\quad \bp_i(T_f)=\mathbf g_i,\quad
 \bv_i(0)=\bv_i(T_f)=\mathbf0 .
 \label{eq:endpoint}
\end{equation}
Both endpoints are inside the free workspace and mutually separated. In the following, $\bv_i^0=\bv_i^f=\mathbf0$ denote these prescribed endpoint velocities. Vector inequalities are understood componentwise.

\subsection{Motion Model and Reduced Variables}
We consider hovering-capable rotary-wing UAVs. Their planning-level translational dynamics are
\begin{equation}
 \dot{\bp}_i=\bv_i,\qquad \dot{\bv}_i=\ba_i.
 \label{eq:continuous}
\end{equation}
Let $h=T_f/K$ and $t_k=kh$, $k=0,\ldots,K$, and take acceleration to be constant on each interval. Exact integration yields
\begin{subequations}\label{eq:discrete}\begin{align}
 \bv_{i,k+1}&=\bv_{i,k}+h\ba_{i,k},\label{eq:velocityupdate}\\
 \bp_{i,k+1}&=\bp_{i,k}+\frac h2(\bv_{i,k}+\bv_{i,k+1}).\label{eq:positionupdate}
\end{align}\end{subequations}
Eliminating acceleration reduces the physical decision fields to $\mathbf P,\mathbf V$:
\begin{equation}
 \ba_{i,k}=\frac{\bv_{i,k+1}-\bv_{i,k}}h,\qquad
 \mathbf E_{i,k}=\bp_{i,k+1}-\bp_{i,k}
 -\frac h2(\bv_{i,k+1}+\bv_{i,k}).
 \label{eq:defect}
\end{equation}
An exact discrete trajectory has $\mathbf E_{i,k}=\mathbf0$. During optimization, this relation is relaxed to allow the position and velocity fields to adjust jointly. Section~\ref{sec:arcs} gives the final control-integration procedure.

The Euclidean speed and acceleration limits are $v_{\max}$ and $a_{\max}$. The planning model uses an isotropic translational envelope. Attitude-dependent actuation enters the rigid-body tracking model in Section~\ref{sec:tracking}.
To retain linear constraints while reducing the directional conservatism of an inscribed cube, define
\begin{equation}
 \mathcal P=\operatorname{conv}\!\left(\{\pm\mathbf e_1,\pm\mathbf e_2,\pm\mathbf e_3\}
 \cup\{(\pm1,\pm1,\pm1)/\sqrt3\}\right)
 =\{\mathbf x:U\mathbf x\leq\mathbf b\}.
 \label{eq:polytope}
\end{equation}
Every vertex has unit Euclidean norm. Convexity of the unit ball therefore gives $\mathcal P\subseteq\{\mathbf x:\norm{\mathbf x}_2\leq1\}$. This 14-vertex, 24-face polytope contains the inscribed Cartesian cube and attains the full limit along each coordinate axis. The motion restrictions are
\begin{equation}
 U\bv_{i,k}\leq v_{\max}\mathbf b,\qquad
 U(\bv_{i,k+1}-\bv_{i,k})/h\leq a_{\max}\mathbf b.
 \label{eq:bounds}
\end{equation}
These are sufficient inner approximations of the Euclidean limits. Final acceptance tests the Euclidean speed and acceleration of the reconstructed motion.

\subsection{Workspace and Collision Conditions}
For $\bp\in\R^3$, define the component gap to an obstacle box as
\begin{equation}
 \mathbf d_b(\bp)=\max\{\underline{\mathbf o}_b-\bp,\,
 \bp-\overline{\mathbf o}_b,\,\mathbf0\}.
 \label{eq:gap}
\end{equation}
The physical requirements, imposed over the complete horizon, are
\begin{subequations}\label{eq:avoidance}\begin{align}
 \underline{\mathbf w}+r\one&\leq\bp_i(t)\leq\overline{\mathbf w}-r\one,\\
 \norm{\mathbf d_b(\bp_i(t))}_2&\geq r+c,\qquad \forall b,\\
 \norm{\bp_i(t)-\bp_j(t)}_2&\geq 2r+c,\qquad \forall i<j.
\end{align}\end{subequations}
Here $t\in[0,T_f]$ is physical flight time, $\bp_i(t)$ is the center position of UAV $i$, and $c$ is the required surface-to-surface clearance. The indices $b$ and $i<j$ cover all obstacles and UAV pairs. These conditions apply to the entire motion of the spherical planning footprint.

\subsection{Acceleration-Effort Objective}
With task completion fixed by~\eqref{eq:endpoint}, the nominal objective is
\begin{equation}
 C(\mathbf V)=w_a h\sum_{i=1}^N\sum_{k=0}^{K-1}
 \left\|\frac{\bv_{i,k+1}-\bv_{i,k}}h\right\|_2^2,\qquad w_a>0.
 \label{eq:cost}
\end{equation}
The coefficient $w_a$ weights acceleration effort, and $C(\mathbf V)$ is the weighted sum of squared accelerations integrated over time. This objective promotes smooth velocity changes in the recovered motion. For a dimensionless objective in SI coordinates, $w_a$ has units $\mathrm{s}^3/\mathrm{m}^2$. The flight duration $T_f$ is prescribed independently of computation time. The outer procedure prioritizes an accepted motion within its computation budget. The quadratic objective guides each local solve.

The nominal problem minimizes~\eqref{eq:cost} subject to exact dynamics, fixed boundary states, motion limits, and continuous collision avoidance. Geometric avoidance makes this problem nonconvex. Spatial and relative corridors provide conservative local linear restrictions, and the recovery mechanism handles their compatibility with the motion model.
